\section{Problem Setup}

We study process verification in audited executable domains. Each example has a
problem statement, a reasoning trace, and a marked audited locus indicating the
local step whose validity is scientifically relevant. The core distinction is
between two input settings:
\begin{itemize}[leftmargin=1.5em]
\item \textbf{answer-visible}: the verifier sees the full trace including the
final answer line;
\item \textbf{answer-masked}: the verifier sees the same trace with the answer
line masked.
\end{itemize}
The central question is whether verifier scores reflect local process validity
or whether they are materially distorted by answer-sensitive behavior.

\subsection{Step validity and matched counterfactuals}

Each benchmark item is derived from an executable latent trajectory. We
construct matched quartets with two orthogonal axes:
\begin{itemize}[leftmargin=1.5em]
\item process validity: \texttt{valid} vs.\ \texttt{invalid};
\item answer surface: \texttt{correct} vs.\ \texttt{swapped}.
\end{itemize}
This yields the minimal quartet:
\texttt{valid\_correct}, \texttt{invalid\_correct},
\texttt{valid\_swapped}, and \texttt{invalid\_swapped}. The point of this
construction is that the verifier can be tested under answer-matched local
comparisons. A model that simply prefers traces with a correct-looking answer
will fail this test.

\subsection{Metrics}

We use the following metrics throughout:
\begin{itemize}[leftmargin=1.5em]
\item \auroc: ordinary valid-vs-invalid discrimination over all traces;
\item \amcd: answer-matched counterfactual discrimination, which asks whether a
verifier prefers \texttt{valid} over \texttt{invalid} when the answer surface
is held fixed;
\item \asstotal: average score change under answer-swap interventions;
\item \selectiongain: improvement over random top-of-quartet selection;
\item \exploitability: how often the top-ranked trace is an invalid-but-lucky
one.
\end{itemize}

These metrics serve different purposes. \auroc{} captures global separability,
\amcd{} captures local process grounding, and \asstotal{} measures direct
sensitivity to answer swapping. The downstream metrics test whether these audit
signals matter operationally.
